
\documentclass[11pt,a4paper]{article}
\usepackage[utf8]{inputenc}
\usepackage[T1]{fontenc}
\usepackage[spanish]{babel}
\usepackage{geometry}
\usepackage{graphicx}
\usepackage{array}
\usepackage{tabularx}
\usepackage{makecell}
\usepackage{hyperref}
\usepackage{xcolor}
\usepackage{fancyhdr}
\usepackage{titlesec}
\usepackage{helvet}
\usepackage{ragged2e}
\usepackage{enumitem}
\usepackage{setspace}
\usepackage{amsmath}
\usepackage{booktabs}
\usepackage{float}
\usepackage{listings}
\usepackage{hyperref}
\usepackage{xurl}
\usepackage{microtype}
\usepackage{float}
\emergencystretch=2em
\renewcommand{\familydefault}{\sfdefault}
\geometry{a4paper,left=3cm,right=3cm,top=2.35cm,bottom=2.35cm,headheight=1.9cm,headsep=0.45cm}
\setlength{\parindent}{0pt}
\setlength{\parskip}{0.3em}
\setlength{\tabcolsep}{5pt}
\renewcommand{\arraystretch}{1.13}
\pagestyle{fancy}
\fancyhf{}
\fancyhead[L]{\IfFileExists{reporte_tecnico_overleaf_logo.png}{\includegraphics[height=1.35cm]{reporte_tecnico_overleaf_logo.png}}{\small Universidad Nacional de Loja}}
\fancyhead[R]{\small FEIRNNR -- Carrera de Computación}
\fancyfoot[C]{\thepage}
\renewcommand{\headrulewidth}{0.5pt}
\renewcommand{\footrulewidth}{0pt}
\titleformat{\section}{\bfseries\fontsize{15}{17}\selectfont}{\thesection.}{0.8em}{}
\titlespacing*{\section}{0pt}{1.0ex}{0.3ex}
\newcolumntype{L}[1]{>{\RaggedRight\arraybackslash}p{#1}}
\newcolumntype{Y}{>{\RaggedRight\arraybackslash}X}
\lstset{basicstyle=\ttfamily\footnotesize,breaklines=true,frame=single,showstringspaces=false,columns=fullflexible}
\hypersetup{hidelinks}

\usepackage{calc}
\usepackage{fvextra}
\providecommand{\tightlist}{\setlength{\itemsep}{0pt}\setlength{\parskip}{0pt}}
\begin{document}
\thispagestyle{fancy}
\begin{center}
{\bfseries\fontsize{21}{24}\selectfont Reporte Técnico de Actividades\\ Práctico-Experimentales Nro. 005\par}
\vspace{0.25cm}
{\large Integración, contenedorización y despliegue del backend}
\end{center}

\section{Datos de Identificación del Estudiante y la Práctica}
\begin{tabularx}{\textwidth}{|L{0.43\textwidth}|Y|}
\hline
\textbf{Nombre del estudiante(s)} & [PENDIENTE: integrantes reales] \\ \hline
\textbf{Asignatura} & Desarrollo Basado en Plataformas (referencia APE4; confirmar) \\ \hline
\textbf{Ciclo} & 5 A (referencia APE4) \\ \hline
\textbf{Unidad} & 1 (referencia APE4) \\ \hline
\textbf{Resultado de aprendizaje} & No aportado específicamente en la guía resumida de APE5. \\ \hline
\textbf{Práctica Nro.} & 005 \\ \hline
\textbf{Título de la Práctica} & Integración, contenedorización y despliegue del backend \\ \hline
\textbf{Nombre del Docente} & Edison Leonardo Coronel Romero (referencia APE4; confirmar) \\ \hline
\textbf{Fecha} & Semana 5; [PENDIENTE: fecha real] \\ \hline
\textbf{Horario} & No especificado en la guía resumida; pendiente. \\ \hline
\textbf{Lugar} & No especificado en la guía resumida; pendiente. \\ \hline
\textbf{Tiempo planificado en el Sílabo} & 3 horas \\ \hline
\end{tabularx}
\textbf{Estado:} preparación documental y artefactos; ejecución funcional, corrección y despliegue pendientes.\par

\section{Objetivo(s) de la
Práctica}\label{APE5-semana-5-objetivos-de-la-pruxe1ctica}

Integrar, contenedorizar y desplegar el backend; publicar documentación
de API y evidencias de pruebas; mantener producto versionado,
instrucciones reproducibles y contribución explicada. La guía se recibió
resumida, por lo que esta organización desarrolla sus puntos sin
inventar un procedimiento oficial numerado ni preguntas de control no
aportadas.

\section{Materiales, Reactivos, Equipos y
Herramientas}\label{APE5-semana-5-materiales-reactivos-equipos-y-herramientas}

Repositorio, especificación/arquitectura/contratos/tests anteriores;
Java/Spring/PostgreSQL del producto; modelo de amenazas y riesgos
Semana4. Se elaboran Dockerfile, Compose, contexto restringido y ejemplo
env. Motor Docker/Podman no disponible, imágenes no descargadas, no
build ni despliegue. Seguridad/DRF/sobre y aceptación PostgreSQL siguen
abiertos. No hay proveedor, credencial de publicación ni URL desplegada
aportados.

\section{Procedimiento / Metodología
Ejecutada}\label{APE5-semana-5-procedimiento--metodologuxeda-ejecutada}

\begin{enumerate}
\def\labelenumi{\arabic{enumi}.}
\tightlist
\item
  Revisar correspondencia de flujos/requisitos y límites de evidencia
  APE3/4.
\item
  Seleccionar empaquetado Spring \textbf{solo como propuesta para el
  backend existente}, sin resolver unilateralmente DRF ni adaptar
  cliente/contrato.
\item
  Preparar Dockerfile con imagen runtime por argumento, JAR explícito,
  usuario no root y puerto8080; el contexto admite solo app.jar y no
  incluye .git ni env del host.
\item
  Preparar Compose con DB PostgreSQL18 sin puerto público, red interna,
  volumen, healthcheck, límites y API loopback18095; secrets fuera de
  Git, prod/demo/cookie explícitos. Las imágenes deben fijarse por
  digest y existir localmente; no se afirma esa disponibilidad.
\item
  Documentar construcción offline/aislada, hash/commit, migraciones,
  fixture, TLS/cookie y permisos efectivos, backup/restauración y
  rollback de imagen/schema; nada se ejecuta por escribir la receta.
\item
  Referenciar documentación de contrato existente y plan de publicación
  estática separada. Prod deshabilita Swagger por configuración: no se
  activa automáticamente para satisfacer ``publicar documentación''.
\item
  Preparar registro de estado de despliegue y checklist de aceptación;
  sin URL ni logs de arranque nuevos. No se usa la disponibilidad local
  temporal de APE3 como despliegue APE5.
\end{enumerate}

\url{INSTRUCCIONES_DESPLIEGUE.md} es la receta completa, con condiciones
y comandos previstos. El motor debe validar Compose y permisos; revisar
sintaxis YAML con una librería no equivale a
\path|docker compose config|, build o aislamiento. Esta entrega mantiene
la constitución: Docker no es único camino obligatorio y Actions no se
dispara automáticamente.

La propuesta usa propietarioDB para migración/arranque por simplicidad
del fixture y lo declara: no alcanza mínimo privilegio productivo. Se
requiere separar migrador/runtime y probar permisos antes de publicar.
Cookie Secure=true exige TLS/proxy para un recorrido de navegador: no
desactivarla en prod para una demo. No se incluyó un proxy/certificado
ficticio ni se afirma CORS configurado.

\section{Resultados}\label{APE5-semana-5-resultados}

\textbf{Producto requerido }: Integración funcional

\textbf{Evidencia disponible }: Código y resultados históricos APE3,
matrices

\textbf{Estado }: Parcial; no nueva integración ejecutada.

\medskip\hrule\medskip

\textbf{Producto requerido }: Contenedorización

\textbf{Evidencia disponible }: Dockerfile y Compose versionables,
contexto allowlist

\textbf{Estado }: Artefactos preparados; no construidos/arrancados.

\medskip\hrule\medskip

\textbf{Producto requerido }: Despliegue / URL o evidencia

\textbf{Evidencia disponible }: estado-despliegue.json declara ausencia

\textbf{Estado }: Pendiente; localhost previsto no es URL desplegada.

\medskip\hrule\medskip

\textbf{Producto requerido }: Instrucciones reproducibles

\textbf{Evidencia disponible }: Receta con imágenes/digests, env,
DB/files, pruebas y rollback

\textbf{Estado }: Documentadas, sin reproducción real certificada.

\medskip\hrule\medskip

\textbf{Producto requerido }: API y pruebas publicadas

\textbf{Evidencia disponible }: Contratos locales y documentación
existente; no hosting

\textbf{Estado }: Disponibles localmente, publicación pendiente.

\medskip\hrule\medskip

\textbf{Producto requerido }: Repositorio versionado/aporte

\textbf{Evidencia disponible }: Archivos y manifiesto; no commit/PR

\textbf{Estado }: Contribución real por completar y versionar.

\medskip\hrule\medskip

\textbf{Producto requerido }: Seguridad básica

\textbf{Evidencia disponible }: Configuración propuesta, controles
previos y SEC pendientes

\textbf{Estado }: No certificada ni cerrada.

\medskip\hrule\medskip

\section{Preguntas de Control}\label{APE5-semana-5-preguntas-de-control}

La guía resumida de APE5 no aportó preguntas numeradas. No se inventan
preguntas oficiales. Para revisar la reproducción: qué commit/JAR/image
se desplegó, qué env/permisos fueron efectivos, qué flujo/pruebas
pasaron, cómo persisten/restauran datos, y qué cambió cada integrante.
Son preguntas de apoyo, no enunciados del docente.

\section{Conclusiones}\label{APE5-semana-5-conclusiones}

Se preparó un paquete de despliegue que evita enviar secretos al
contexto y explicita dependencias/configuración. \textbf{La práctica no
está completamente realizada:} faltan build, despliegue, pruebas,
publicación, correcciones de seguridad y evidencia individual
versionada. Tener un Dockerfile no equivale a un backend desplegado.

\section{Recomendaciones}\label{APE5-semana-5-recomendaciones}

Resolver tecnología/contrato; verificar controles pendientes; disponer
de sandbox/motor e imágenes auditadas; reconstruir artefacto del commit
actual; ejecutar pruebas y migraciones PostgreSQL; validar TLS/mínimo
privilegio y restauración; publicar solo documentación/servicio
autorizados con evidencia concreta. No limpiar volúmenes ni enviar
secretos a logs/repo.

\section{Bibliografía /
Referencias}\label{APE5-semana-5-bibliografuxeda--referencias}

\href{../soporte-semanas-4-5/FUENTES.md}{Fuentes compartidas}, D-01--04
y F-01--12. Docker Compose/redes y PostgreSQL oficial respaldan la forma
propuesta, no su ejecución. Documento Boot consultado4.1.2 no actualiza
la dependencia configurada4.1.1.

\nocite{*}
\bibliographystyle{IEEEtran}
\bibliography{referencias}
\section{Anexos}\label{APE5-semana-5-anexos}

\path|despliegue/Dockerfile|, \path|compose.yaml|, \path|.env.example|,
\path|.gitignore|, \path|context/.dockerignore|, instrucciones, estado
de despliegue, trazabilidad/bitácora AA5 y nota individual. No se
entrega un JAR antiguo, una imagen no construida como si existiese ni
una URL de producción inventada.






\subsection{Fuente propuesta: Dockerfile}
\begin{Verbatim}[breaklines=true,breakanywhere=true,fontsize=\footnotesize]
# Proposal only: not built or deployed. Supply an audited locally available image.
ARG RUNTIME_IMAGE
FROM ${RUNTIME_IMAGE}
WORKDIR /app
RUN mkdir -p /data/files && chown -R 10001:10001 /data
COPY --chown=10001:10001 app.jar /app/app.jar
USER 10001:10001
EXPOSE 8080
ENTRYPOINT ["java", "-XX:MaxRAMPercentage=65", "-jar", "/app/app.jar"]

\end{Verbatim}

\subsection{Fuente propuesta: compose.yaml}
\begin{Verbatim}[breaklines=true,breakanywhere=true,fontsize=\footnotesize]
name: sigeac-academic-local
services:
  db:
    image: ${POSTGRES_IMAGE:?Set an audited local PostgreSQL 18 image pinned by digest}
    pull_policy: never
    environment:
      POSTGRES_DB: sec_unl
      POSTGRES_USER: sec_owner
      POSTGRES_PASSWORD: ${SEC_DB_PASSWORD:?Set outside version control}
    volumes:
      - pgdata:/var/lib/postgresql
    networks: [backend]
    healthcheck:
      test: [CMD-SHELL, "pg_isready -U sec_owner -d sec_unl"]
      interval: 5s
      timeout: 3s
      retries: 12
    mem_limit: 512m
    cpus: 1
    pids_limit: 100
    security_opt: [no-new-privileges:true]
  api:
    image: sigeac-academic-local:unverified
    pull_policy: never
    build:
      context: ./context
      dockerfile: ../Dockerfile
      network: none
      args:
        RUNTIME_IMAGE: ${JAVA_RUNTIME_IMAGE:?Set an audited local Java 25 JRE image pinned by digest}
    environment:
      SPRING_PROFILES_ACTIVE: prod
      SERVER_ADDRESS: 0.0.0.0
      SEC_DB_URL: jdbc:postgresql://db:5432/sec_unl
      SEC_DB_USER: sec_owner
      SEC_DB_PASSWORD: ${SEC_DB_PASSWORD:?Set outside version control}
      SEC_AUTH_HMAC_SECRET: ${SEC_AUTH_HMAC_SECRET:?Set a high entropy secret outside Git}
      SEC_DEMO_DATA: "false"
      SEC_COOKIE_SECURE: "true"
      SEC_FILE_ROOT: /data/files
      SPRINGDOC_API_DOCS_ENABLED: "false"
      SPRINGDOC_SWAGGER_UI_ENABLED: "false"
    depends_on:
      db:
        condition: service_healthy
    ports:
      - "127.0.0.1:18095:8080"
    volumes:
      - files:/data/files
    networks: [backend]
    read_only: true
    tmpfs:
      - /tmp:size=64m,mode=1777
    cap_drop: [ALL]
    security_opt: [no-new-privileges:true]
    mem_limit: 768m
    cpus: 1
    pids_limit: 100
networks:
  backend:
    internal: true
volumes:
  pgdata:
  files:

\end{Verbatim}
\section{Declaración de uso de Inteligencia
Artificial}\label{APE5-semana-5-declaraciuxf3n-de-uso-de-inteligencia-artificial}

Se utilizó ChatGPT/Codex para inspeccionar fuentes y preparar
documentos/configuración. Las comprobaciones de artefactos son estáticas
y no validan seguridad efectiva. La contribución/aceptación humanas y la
ejecución de build/despliegue siguen pendientes; no se atribuyen a
integrantes desconocidos.

\end{document}
